% from ega4-1-fr.tex
\oldpage[IV-1]{239}

et soit $f:X'\to X$ le morphisme structural~; on vient de voir que
l'immersion canonique $Z'\to X$ est de présentation finie, et il en est de
même de l'immersion ouverte $U\to X$, qui est quasi-compacte puisque $U$ et
$X$ sont affines~; on conclut donc de (1.6.2, (iv)) que $f$ est de
présentation finie, et l'on a $f(X')=Z'\cap U$ (I, 3.4.8), ce qui achève la
démonstration.
\end{proof}

\begin{theorem}[1.8.4]
\label{IV.1.8.4-fr}
\textup{(Chevalley).} --- Soit $f:X\to Y$ un morphisme de présentation
finie. Pour toute partie localement constructible $Z$ de $X$, $f(Z)$ est
localement constructible dans $Y$.
\end{theorem}

\begin{proof}
Soient $y\in Y$ et $V$ un voisinage ouvert affine de $y$~; comme $f$ est
quasi-compact et quasi-séparé, il en est de même de sa restriction
$f^{-1}(V)\to V$, donc $f^{-1}(V)$ est un préschéma quasi-compact et
quasi-séparé (1.2.3)~; comme $Z\cap f^{-1}(V)$ est constructible dans
$f^{-1}(V)$ (1.8.1), on voit qu'on peut se borner au cas où $Y$ est affine
et $Z$ constructible~; $X$ est alors quasi-compact et quasi-séparé, donc
(1.8.3), il existe un morphisme de présentation finie $g:X'\to X$ tel que
$Z=g(X')$~; comme $f\circ g$ est de présentation finie, on voit qu'on est
ramené au cas où $Z=X$~; autrement dit, il suffira de prouver le
\end{proof}

\begin{lemma}[1.8.4.1]
\label{IV.1.8.4.1-fr}
Soient $Y$ un schéma affine, $f:X\to Y$ un morphisme quasi-compact et
localement de présentation finie~; alors $f(X)$ est une partie constructible
de $Y$.
\end{lemma}

\begin{proof}
Comme $X$ est quasi-compact, il est réunion finie d'ouverts affines~; on peut
donc se borner au cas où $Y=\operatorname{Spec}(A)$,
$X=\operatorname{Spec}(B)$, $B$ étant une $A$-algèbre de présentation finie
(1.4.6). Or, on a le
\end{proof}

\begin{lemma}[1.8.4.2]
\label{IV.1.8.4.2-fr}
Soient $A_0$ un anneau, $(A_\lambda)_{\lambda\in L}$ un système inductif de
$A_0$-algèbres, $A=\varinjlim A_\lambda$~; si $B$ est une $A$-algèbre de
présentation finie, il existe un $\lambda$ et une $A_\lambda$-algèbre de
présentation finie $B_\lambda$ tels que $B$ soit isomorphe à
$B_\lambda\otimes_{A_\lambda}A$.
\end{lemma}

\begin{proof}
Par hypothèse, $B$ est isomorphe à une algèbre de la forme
$A[T_1,\ldots,T_n]/\mathfrak J$, où les $T_i$ sont des indéterminées et
$\mathfrak J$ un idéal de type fini~; soit $(F_j)$ ($1\leq j\leq m$) un
système de générateurs de $\mathfrak J$. Soit
$\varphi_\lambda:A_\lambda\to A$ l'homomorphisme canonique. Comme $L$ est
filtrant, il existe $\lambda$ tel que chacun des coefficients de chacun des
polynômes $F_j$ soit l'image par $\varphi_\lambda$ d'un élément de
$A_\lambda$. Autrement dit, il existe un système de polynômes
$(F_{j\lambda})_{1\leq j\leq m}$ de $A_\lambda[T_1,\ldots,T_n]$ tels que
$F_j=\varphi_\lambda(F_{j\lambda})$ pour $1\leq j\leq m$. Si
$\mathfrak J_\lambda$ est l'idéal de $A_\lambda[T_1,\ldots,T_n]$ engendré
par les $F_{j\lambda}$, $\mathfrak J$ est l'image de
$\mathfrak J_\lambda\otimes_{A_\lambda}A$ dans
$A[T_1,\ldots,T_n]=A_\lambda[T_1,\ldots,T_n]\otimes_{A_\lambda}A$~;
l'anneau $B_\lambda=A_\lambda[T_1,\ldots,T_n]/\mathfrak J_\lambda$ répond à
la question.
\end{proof}

On appliquera ici ce lemme en considérant $A$ comme limite inductive de ses
sous-$\mathbf Z$-algèbres \emph{de type fini}. On voit donc qu'il existe une
telle sous-$\mathbf Z$-algèbre $A_0$ de $A$, et une $A_0$-algèbre de
présentation finie $B_0$ telles que $B$ soit isomorphe à
$B_0\otimes_{A_0}A$~; posons $X_0=\operatorname{Spec}(B_0)$,
$Y_0=\operatorname{Spec}(A_0)$, de sorte que l'on a $X=X_0\times_{Y_0}Y$,
$f$ étant la projection $X\to Y$~; soient $f_0:X_0\to Y_0$,
$g_0:Y\to Y_0$ les morphismes structuraux~; il résulte de (I, 3.4.8) que
l'on a $f(X)=g_0^{-1}(f_0(X_0))$~; compte tenu de (1.8.2), il suffit donc de
montrer que $f_0(X_0)$ est constructible. Autrement dit, on est finalement
ramené à démontrer le

\begin{corollary}[1.8.5]
\label{IV.1.8.5-fr}
Soient $Y$ un préschéma noethérien, $f:X\to Y$ un morphisme de type fini.
Pour toute partie constructible $Z$ de $X$, $f(Z)$ est une partie
constructible de $Y$.
\end{corollary}

